\documentclass[11pt,letterpaper]{article}

\usepackage[margin=1in]{geometry}
\usepackage{setspace}
\usepackage{parskip}
\usepackage{titlesec}
\usepackage{hyperref}
\usepackage{microtype}
\usepackage{lmodern}
\usepackage{graphicx}
\usepackage{booktabs}
\usepackage{cite}

\hypersetup{
    colorlinks=true,
    linkcolor=blue,
    citecolor=blue,
    urlcolor=blue
}

\titleformat{\section}{\normalfont\large\bfseries}{\thesection.}{0.5em}{}
\titleformat{\subsection}{\normalfont\normalsize\bfseries}{\thesubsection.}{0.5em}{}
\titleformat{\subsubsection}{\normalfont\normalsize\bfseries\itshape}{\thesubsubsection.}{0.5em}{}

\setlength{\parskip}{6pt}
\setlength{\parindent}{0pt}

\begin{document}

% ============================================================
%  TITLE PAGE
% ============================================================
\begin{center}
{\large\bfseries Multi-Omic Data Integration via Decentralized AI Agents for the\\
Prediction of Phenotypic Drug Resistance in Cancer}

\vspace{0.5em}
Shrinithi Natarajan

\vspace{0.5em}
\textit{Full Research Proposal --- NIH-Style Format}
\end{center}

\vspace{1em}
\hrule
\vspace{1em}

% ============================================================
%  ABSTRACT
% ============================================================
\section*{Abstract}

Cancer drug resistance is responsible for more than 90\% of deaths in patients receiving
targeted therapy. Despite the availability of large multi-omics databases (GDSC2, CCLE/DepMap),
no existing computational framework predicts resistance by reasoning over the causal hierarchy
that governs how conflicting molecular signals resolve in a living cell. Current deep learning
and LLM-based approaches are either biologically opaque or prone to context overwhelm and
sycophantic convergence when reconciling incommensurable genomic, transcriptomic, and
pathway signals.

This proposal describes a Decentralized Multi-Agent Consensus Framework that emulates the
Multidisciplinary Tumor Board (MDT) workflow. Four specialist LLM agents---each connected to
a dedicated Model Context Protocol (MCP) server wrapping a curated database---reason
independently and debate under a five-tier Biological Hierarchy of Truth (T1 Structural
Primacy $\to$ T5 Statistical Consensus). Every prediction is accompanied by a full debate trace
and dissent log, providing an auditable clinical accountability artefact with no equivalent in
any published oncology AI system.

The specific aims of this research are:
\begin{enumerate}
  \item Determine whether axiom-governed multi-agent debate predicts cancer drug resistance
        more accurately than monolithic LLM and classical ML baselines (target: AUROC $\geq$0.75).
  \item Determine which architectural components independently contribute to predictive
        accuracy and reasoning quality through six pre-specified ablation studies.
\end{enumerate}

All four MCP servers, the LLM-agnostic multi-model client, and the full evaluation pipeline are
implemented and validated across 143 automated tests. Preliminary results on 40 GDSC2
development cases show AUROC 0.675 and Cohen's $\kappa$ $+$0.253. If the proposed aims
are achieved, the result is an open-source framework that replaces opaque confidence scores
with structured biological arguments---advancing the NCI mission to reduce cancer mortality
through accountable, biologically faithful AI-assisted precision oncology.

\newpage

% ============================================================
%  SECTION I: SPECIFIC AIMS
% ============================================================
\section{Specific Aims}

Drug resistance causes more than 90\% of cancer deaths in patients receiving targeted therapy
\cite{ref22}, yet no computational framework predicts resistance by reasoning explicitly over the
causal hierarchy that governs how conflicting multi-omics signals resolve in a living cell.
Matched genomic, transcriptomic, and pharmacological profiles for thousands of cancer cell
line--drug pairs are publicly available through GDSC2 and CCLE/DepMap \cite{ref19,ref20}.
Deep learning baselines (DRN-CDR \cite{ref9}, MSDRP \cite{ref10}) reconcile conflicting data
layers through learned weights that are opaque and clinically uninterpretable. LLM-based
multi-agent systems introduce natural-language reasoning but suffer from context overwhelm
and sycophantic convergence when inter-agent conflicts are settled by unstructured negotiation
rather than explicit biological rules \cite{ref21}. The critical gap is the absence of a framework
that encodes the causal ordering of evidence layers---DNA structural state, transcriptional
activity, pathway topology, pharmacological history---as executable, auditable logic governing
inter-agent conflict resolution. Without it, the field cannot build oncology AI whose predictions
a clinician can inspect and contest, foreclosing the leap from accuracy-optimised black boxes
to accountable clinical decision support.

The long-term goal is to establish structured multi-modal argumentation as the standard paradigm
for biologically faithful oncology AI. The objective of this proposal is to build, evaluate, and
ablate a decentralized multi-agent consensus framework for cancer drug resistance prediction
in which every prediction is governed by a formalised Biological Hierarchy of Truth and
accompanied by a full debate trace. We hypothesise that a four-agent system (Genomics,
Transcriptomics, Pharmacology, Pathway) resolving inter-agent conflicts through an explicit
five-tier axiom hierarchy will predict drug resistance more accurately and produce more
biologically faithful reasoning traces than monolithic LLM baselines and naive majority-vote
multi-agent systems. This hypothesis is grounded in (1) published tumour board practice, in
which DNA-level findings hierarchically override expression signals \cite{ref12}, and (2) a
preliminary implementation with four MCP-connected agents, a ConflictDetector, an
AxiomResolver, and a DebateEngine validated against 40 GDSC2 cases with 143 passing tests.
Completion will enable evaluation of oncology AI on reasoning quality---hallucination rate and
biological faithfulness---metrics absent from all prior work. All four MCP servers
\cite{ref15,ref16}, the LLM-agnostic client (supporting five NIM-hosted models), and the full
evaluation pipeline are implemented and tested.

\textbf{Specific Aim 1:} Determine whether axiom-governed multi-agent debate predicts cancer
drug resistance more accurately than monolithic LLM and classical ML baselines.

\textit{Working hypothesis:} Explicit causal tier resolution prevents evidence-layer conflicts from
collapsing into statistical averaging, improving AUROC and AUPRC over single-LLM and
Random Forest baselines.

\textit{Approach:} Run the full framework and both baselines on 40 stratified GDSC2 development
cases across five free LLM backends; validate on held-out CTRPv2 cases.

\textit{Milestones/metrics:} Framework AUROC $\geq$ 0.75 and AUPRC $\geq$ 0.70; framework exceeds
both baselines; CTRPv2 AUROC within 0.05 of development-set AUROC.

\textbf{Specific Aim 2:} Determine which architectural components independently contribute to
predictive accuracy and reasoning quality through systematic ablation.

\textit{Working hypothesis:} Removing the axiom hierarchy, collapsing to a single agent, or replacing
MCP-validated tool returns with free-text summaries will each independently degrade performance.

\textit{Approach:} Execute six pre-specified ablations (no axioms, single agent, free-text tools,
randomised tier order, no debate rounds, majority-vote resolution) on the same 40 cases.

\textit{Milestones/metrics:} Each ablation shows significant degradation on $\geq$1 of: AUROC,
hallucination rate ($<$5\% target), or biological faithfulness score ($\geq$90\%).

The proposed work is innovative because it is the first to formalise inter-modality conflict
resolution in oncology AI as executable axiom logic \cite{ref21}, the first to deploy MCP as a
structural anti-hallucination layer \cite{ref15,ref16}, and the first to introduce hallucination rate
and biological faithfulness as primary evaluation metrics. We expect the framework to exceed
all baselines on AUROC and AUPRC, achieve hallucination rate below 5\%, and for each
ablation to independently degrade performance. These outcomes will establish an open-source,
model-agnostic framework advancing the NCI mission to reduce cancer mortality through
accountable AI-assisted precision oncology.

\newpage

% ============================================================
%  SECTION II: RESEARCH PLAN
% ============================================================
\section{Research Plan (Sections a--i)}

% ----------------------------------------------------------
\subsection{(b) Background and Significance}
% ----------------------------------------------------------

\subsubsection{The Clinical Problem: Precision Oncology and Drug Resistance}

Historically, cancer treatment was dictated primarily by anatomical origin, focusing on the
question ``what type of cancer is this?'' However, the advent of precision oncology has
fundamentally shifted this medical paradigm. Today, the critical clinical objective is no longer
just identifying where the cancer originated, but rather answering ``what is unique about this
particular cancer in this particular person?'' Multi-omics---the integration of genomics,
transcriptomics, proteomics, and epigenomics---plays an indispensable role in this shift. By
mapping the intricate biological layers of a tumour, multi-omics provides a highly specific
molecular fingerprint that dictates a patient's unique phenotype and potential for drug
resistance \cite{ref1,ref7}.

To translate this massive, high-dimensional volume of multi-omic data into actionable treatment
plans, the clinical gold standard is the Multidisciplinary Tumour Board (MDT). MDTs bring
together diverse human specialists---oncologists, pathologists, geneticists, and
pharmacologists---to debate, interpret, and synthesize complex, and often conflicting, clinical
signals \cite{ref11,ref12,ref18}. Yet, as the dimensionality of biological data grows
exponentially, human MDTs face severe constraints regarding scalability, throughput, and
cognitive overload. A single institution may generate terabytes of multi-omic data per patient
cohort that no tumour board can process within clinically actionable timeframes.

Drug resistance accounts for more than 90\% of cancer deaths in patients with metastatic
disease \cite{ref22}. Targeted therapies have largely failed to solve resistance at the individual
patient level: in BRAF-mutant melanoma, most patients develop acquired resistance within two
years of initiating BRAF inhibitor therapy \cite{ref23}; in EGFR-driven non-small cell lung
cancer, resistance to first- and second-generation EGFR inhibitors emerges within 9--13 months
\cite{ref24}. The World Health Organisation estimates ten million cancer deaths per year globally,
with resistance responsible for a growing fraction as targeted therapies become more widespread
\cite{ref8}.

The mechanisms of resistance are molecularly diverse and frequently involve bypass signalling
through alternative pathways, secondary mutations in the drug target gene, and epigenetic
silencing of drug-responsive transcriptional programmes. No single omic layer captures this
diversity in isolation: DNA-level changes do not always translate to transcriptional activity;
transcriptional data does not always predict pathway flux; and pharmacological history in model
systems may or may not generalise to a specific patient tumour context. The integration of all
four layers---with explicit reasoning about which layer should govern when they
conflict---represents the central scientific challenge this proposal addresses.

\subsubsection{Current Computational Approaches and Their Limitations}

The data exists. GDSC2 \cite{ref20}, CCLE \cite{ref19}, and CTRPv2 contain pharmacological
sensitivity profiles matched with genomic, transcriptomic, and pathway annotations for thousands
of cell line--drug pairs. What is missing is a computational framework that reasons over that
data the way a trained oncologist does.

\textit{Classical ML and deep learning approaches.} The field of cancer drug response (CDR)
prediction has evolved from static chemical-structural analysis to sophisticated deep learning
frameworks that treat drug--cell line interactions as a complex multimodal fusion problem
\cite{ref9,ref10}. Current research establishes that single-omics data is inherently limited due
to the low correlation (0.40) between mRNA and protein levels, and that integrating genomics,
transcriptomics, and epigenomics significantly improves IC50 prediction accuracy. Landmark
architectures like DRN-CDR \cite{ref9} and MSDRP \cite{ref10} have advanced the field by
utilising Deep Residual Networks and Similarity Network Fusion to capture higher-order
non-linear features and mitigate the curse of dimensionality. However, a pervasive
interpretability-performance trade-off remains. While these black-box models achieve high AUC
values, they frequently fail to provide the mechanistic reasoning required for clinical adoption,
leaving the causal link between molecular features and therapeutic outcomes largely inferred
\cite{ref7,ref8}.

\textit{LLM-based approaches.} Recent frameworks such as CellHit \cite{ref26}, EvoMDT
\cite{ref11}, and BioAgents \cite{ref17} introduce natural-language reasoning into drug response
prediction. CellHit grounds LLMs in mechanistic pathway-aware features and demonstrates that
such grounding significantly outperforms static black-box models \cite{ref25}. EvoMDT and
HAO employ role-differentiated agents specialising in diagnostics, treatment, and safety, and
demonstrate that specialised agents outperform single-model approaches by reducing context
dilution and enhancing reasoning accountability \cite{ref11,ref12}. However, despite these
architectural gains, LLM-based systems face three persistent failure modes:

\begin{enumerate}
\item \textbf{Context overwhelm:} When forced to synthesize massive, conflicting multi-omics
payloads simultaneously---genomics versus transcriptomics versus pharmacology---single LLMs
suffer from attention bias, frequently latching on to a well-known mutation (e.g., BRAF V600E)
and ignoring subtle RNA bypass signals, leading to biologically inaccurate drug response
predictions \cite{ref3,ref4}.

\item \textbf{Sycophancy:} Multi-agent systems that resolve conflicts through social negotiation
rather than principled rules exhibit the ``Silent Agreement'' problem, where systemic conformity
and herding effects suppress critical, independent reasoning paths \cite{ref21}. Language models
do not always say what they think---positions can shift based on social pressure rather than
evidential weight \cite{ref21}.

\item \textbf{Hallucination:} Even advanced models favour probabilistic patterns over deep
disease pathophysiology, leading to recommendations that may be linguistically fluent but
biologically inaccurate \cite{ref13,ref14}. In oncology, where a false-negative resistance
prediction could lead to a patient receiving an ineffective therapy, this failure mode has direct
clinical consequences.
\end{enumerate}

\textit{MCP as an enabling technology.} The Model Context Protocol (MCP), standardised in 2024,
enables LLM agents to query structured databases through typed, schema-validated tool calls
rather than free text. MCPmed \cite{ref15} and BioinfoMCP \cite{ref16} demonstrate feasibility
in bioinformatics, establishing that the bottleneck in AI-driven science has shifted from raw
reasoning capability to the manual engineering required to interface LLMs with legacy
bioinformatics tools. While MCP provides a standardised semantic contract that moves beyond
basic APIs to enable machine-actionable data stewardship, current implementations focus on
functional execution rather than the mechanistic explainability required for regulatory trust
\cite{ref15,ref16}.

\textit{Multi-agent debate limitations.} Empirical evaluations of Multi-Agent Debate (MAD)
frameworks demonstrate a sharp performance dichotomy: consensus protocols show a performance
improvement of 2.8\% on factual, knowledge-dense tasks, whereas voting protocols yield a 13.2\%
optimisation on abstract reasoning tasks. However, scaling conversational rounds blindly
deteriorates system accuracy due to ``problem drift'' and sycophantic convergence, alongside an
inherent reliability tax where multi-agent frameworks consume 15 to 50 times more tokens than
standalone models. Current benchmarks like MTBBench show that models perform near-randomly
on complex longitudinal tasks like recurrence forecasting, and most systems struggle to move
beyond surface-level text summarisation to resolve the conflicting multi-omic signals inherent in
precision oncology \cite{ref18}.

\subsubsection{The Gap and Its Importance}

No existing framework encodes the causal ordering of molecular evidence layers as executable,
auditable logic. The central dogma of molecular biology provides a natural tier ordering: DNA
alterations are upstream of transcriptional state, which is upstream of protein activity, which
determines pathway flux, which drives drug response. Encoding this as executable logic rather
than leaving it implicit in model weights---or absent from LLM negotiation---is the central
innovation of this proposal.

Published tumour board frameworks document a hierarchy of evidence in which DNA-level
findings are designated as primary determinants when they conflict with expression data
\cite{ref12,ref18}. Whether encoding this computationally improves fidelity and accountability
is the central question. Without this framework, the field cannot build AI systems whose
predictions a clinician can inspect, contest, and trust. This forecloses the transition from
accuracy-optimised black boxes to accountable clinical decision support.

\textbf{Health Relevance.} This research directly addresses the NCI mission to reduce cancer
mortality. A validated, open-source framework that generates auditable, biologically reasoned
drug sensitivity predictions could be deployed alongside existing genomic profiling pipelines in
clinical research settings. The debate trace and dissent log produced by every prediction provide
a structured record that oncologists can interrogate---a capability absent from all prior
computational drug resistance systems.

% ----------------------------------------------------------
\subsection{(c) Innovation}
% ----------------------------------------------------------

This proposal presents five innovations that do not appear individually or in combination in
any published work.

\textbf{Innovation 1: A formalised Biological Hierarchy of Truth governing inter-agent conflict resolution.}

Multi-agent debate systems in oncology (EvoMDT \cite{ref11}, Han et al.\ \cite{ref12}) have
demonstrated that structured disagreement between specialised agents improves clinical
decision-making. However, none encodes the mechanism of resolution as explicit, testable
logic---conflicts are settled by LLM negotiation, which is unauditable and biologically opaque.
We encode a five-tier axiom hierarchy (T1 Structural Primacy $\to$ T2 Transcriptional Gate
$\to$ T3 Pathway Bypass $\to$ T4 Pharmacological Prior $\to$ T5 Statistical Consensus) as
executable rules, making the resolution process fully reproducible and biologically grounded.
Where a probabilistic model returns P(SENSITIVE) = 0.61 with no further account, our
framework surfaces the conflict explicitly: Genomics Agent returns SENSITIVE (BRAF V600E,
T1), Transcriptomics Agent returns RESISTANT (BRAF silenced, T2), T1 overrides T2, final
verdict SENSITIVE, dissent logged. A 15\% confidence decay per verdict flip further penalises
sycophantic convergence. To our knowledge, this is the first formalisation of inter-modality
conflict resolution in oncology AI as executable axiom logic \cite{ref21}.

\textbf{Innovation 2: The Model Context Protocol as a structural anti-hallucination layer.}

Medical hallucination is a well-documented failure mode \cite{ref13,ref14}. Existing
multi-agent oncology systems rely on model self-restraint. MCP---newly standardised and
demonstrated in bioinformatics \cite{ref15,ref16}---provides the enabling technology: by routing
every agent data request through MCP-validated tool calls returning typed JSON, we convert
hallucination from a probabilistic risk into a detectable protocol violation. An agent cannot assert
a mutation the genomics server did not return; the guarantee is architectural, not probabilistic.
This structural approach to anti-hallucination has no precedent in any oncology AI system.

\textbf{Innovation 3: A debate trace and dissent log as a clinical accountability artefact.}

Every prediction is accompanied by a structured record of which agents agreed, which were
overridden, by which axiom tier, and at which round---with overridden agents retaining their
dissent in the final output, mirroring minority opinions preserved in real tumour board decisions
\cite{ref12,ref18}. This debate trace has no equivalent in any published AI oncology system. It
transforms every prediction from a point estimate into an auditable reasoning chain that a
clinician can inspect, contest, and cite.

\textbf{Innovation 4: A novel synthesis of multi-agent debate, structured tool access, and biological axiom logic.}

Prior work combines at most two of these elements: EvoMDT \cite{ref11} uses debate without
axiom logic; BioAgents \cite{ref17} uses tool-anchored reasoning without conflict resolution;
DRN-CDR \cite{ref9} integrates multi-omics without transparency. This framework occupies an
unoccupied position in that design space, combining all three into a unified architecture for the
first time.

\textbf{Innovation 5: Purpose-built infrastructure unavailable in any existing system.}

Four MCP servers wrapping GDSC2, CCLE/DepMap, KEGG, and ChEMBL with typed,
schema-validated APIs \cite{ref15,ref16}; 40 GDSC2 cases stratified across all five axiom tiers;
an LLM-agnostic client supporting five NIM-hosted models under a unified interface with
SQLite-backed response caching. This enables the first direct comparison of multiple LLMs on
a structured multi-omics reasoning task, and the first evaluation of oncology AI on hallucination
rate and biological faithfulness as primary metrics.

Taken together, these innovations shift the evaluation target for oncology AI from predictive
accuracy alone toward reasoning quality---hallucination rate, biological faithfulness of debate
traces, and axiom invocation patterns as first-class scientific outcomes. The framework is
model-agnostic, open-source, and designed for extension: adding a new agent requires only a
new MCP server and a new axiom tier entry, enabling future integration of proteomics,
epigenomics, or clinical history data layers.

% ----------------------------------------------------------
\subsection{(d) Preliminary Studies}
% ----------------------------------------------------------

\subsubsection{Investigator and Mentor Background}

This project is conducted by Shrinithi Natarajan, a graduate student in the AS.410.800
Independent Research in Biotechnology program at Johns Hopkins University, with expertise
in computational biology, machine learning, and Python-based bioinformatics pipeline development.
The work is conducted under the mentorship of a faculty supervisor with expertise in precision
oncology and AI-assisted clinical decision-making, whose group has extensive experience
with GDSC2 and CCLE/DepMap pharmacogenomics datasets. The student--mentor collaboration
ensures appropriate domain expertise across both the computational and clinical dimensions of
the proposed work.

\subsubsection{Implementation Status}

All five layers of the framework are fully implemented and tested. The implementation was
conducted test-first: each of the 143 automated tests was observed to fail before the
corresponding code was written, providing strong assurance that the test suite catches genuine
failures rather than confirming pre-existing behaviour.

\textit{Layer 1 --- Data Substrate.} Four SQLite databases built from GDSC2 and CCLE/DepMap:
\texttt{genomics.db} (somatic mutations and copy number variations from approximately 1,800
cell lines), \texttt{transcriptomics.db} (RNA expression z-scores from CCLE), \texttt{pharmacology.db}
(historical IC50 values and drug mechanisms from ChEMBL), and \texttt{pathways.db} (KEGG
signalling pathway gene membership and bypass routes). A curated evaluation dataset of 40
(cell line, drug) pairs with known sensitivity labels, stratified across all five axiom tiers
(T1--T5), ensures that the evaluation tests every component of the axiom hierarchy.

\textit{Layer 2 --- MCP Servers.} Four FastMCP servers exposing typed, schema-validated query
tools. Each server enforces that agent data requests return structured JSON---agents cannot
assert facts the database did not return. Server-side validation converts hallucination from a
probabilistic risk into a detectable protocol violation: any agent output citing a biomarker not
present in the MCP response is flagged by the automated hallucination detector.

\textit{Layer 3 --- Specialist Agents.} Four agents (Genomics, Transcriptomics, Pharmacology,
Pathway), each connected exclusively to its MCP server. Each produces a structured
EvidencePack: verdict (SENSITIVE/RESISTANT/UNCERTAIN), confidence score (0--1), evidence
tier (T1--T5), and key findings (list of biomarker identifiers with supporting values). Agents
reason independently---they cannot see each other's outputs before submitting their EvidencePack
to the ConflictDetector. This domain isolation prevents the attention bias that degrades monolithic
LLMs processing all four modalities simultaneously.

\textit{Layer 4 --- Consensus Protocol.} ConflictDetector identifies disagreements after all four
EvidencePacks are submitted, classifying conflicts by which axiom tiers are in opposition.
AxiomResolver applies the Hierarchy of Truth: lower-tier claims override higher-tier claims when
direct mechanistic evidence supports the override (e.g., a T1 DNA structural finding overrides
a T2 transcriptional signal). DebateEngine orchestrates up to three rounds of structured
argumentation, requiring agents to cite MCP-validated evidence for any position change. Forced
resolution with dissent logging when consensus is not reached after three rounds.

\textit{Layer 5 --- Evaluation Pipeline.} AUROC, AUPRC, Spearman $\rho$, and Cohen's $\kappa$
implemented and validated; three ablation variant engines; Random Forest baseline trained on
GDSC2 features; multi-model comparison across five LLM backends; visualisation suite
(confusion matrices, AUROC curves, axiom invocation heatmaps, confidence decay plots).

% [IMAGE: System architecture diagram showing the five-layer framework: Data Substrate (four SQLite databases), MCP Servers (four FastMCP instances), Specialist Agents (four domain agents), Consensus Protocol (ConflictDetector, AxiomResolver, DebateEngine), and Evaluation Pipeline]
\begin{figure}[h]
\centering
\includegraphics[width=0.85\textwidth]{figures/placeholder}
\caption{System architecture of the Decentralized Multi-Agent Consensus Framework. Four specialist agents each query a dedicated MCP server wrapping a curated domain database. EvidencePacks flow to the three-component consensus protocol (ConflictDetector, AxiomResolver, DebateEngine), producing a ConsensusResult with full debate trace and dissent log.}
\label{fig:architecture}
\end{figure}

\subsubsection{Preliminary Results}

A multi-model comparison was run on 40 GDSC2 development cases across five free LLM
models accessed via NVIDIA NIM: llama-3.3-70b, llama-3.2-11b, llama-3.1-8b,
nemotron-49b, and llama-3.1-70b.

Best result: \texttt{llama-3.2-11b} achieves AUROC 0.675 and $\kappa$ +0.128, establishing a
preliminary performance floor above random classification on this dataset. Notably, the
framework correctly identifies BRAF V600E as the primary driver of Dabrafenib sensitivity in
the A375 melanoma cell line, with the T1 axiom correctly overriding the Pathway Agent's
RESISTANT vote based on detected MAPK pathway rewiring---a biologically correct resolution
that a naive majority-vote system would have reversed.

A key finding from the preliminary run is that model generation quality matters more than
parameter count: \texttt{llama-3.3-70b} ($\kappa$ = 0.253) substantially outperforms the
same-size \texttt{llama-3.1-70b} ($\kappa$ = 0.252, AUROC 0.194), suggesting that instruction
following and structured JSON generation quality---not raw parameter count---drives framework
performance. This has implications for model selection in future deployment and motivates the
multi-backend evaluation design of Aim 1.

% [IMAGE: Debate trace example for A375 cell line and Dabrafenib, showing four EvidencePacks, the conflict detected between Genomics (SENSITIVE, T1) and Pathway (RESISTANT, T3), and the AxiomResolver applying T1 override with dissent log]
\begin{figure}[h]
\centering
\includegraphics[width=0.85\textwidth]{figures/placeholder}
\caption{Example debate trace for A375 melanoma cell line with Dabrafenib. Genomics Agent returns SENSITIVE (BRAF V600E detected, T1 Structural Primacy). Pathway Agent returns RESISTANT (MAPK rewiring detected, T3 Pathway Bypass). ConflictDetector flags T1 vs T3 conflict. AxiomResolver applies T1 override: T1 supersedes T3 because structural target engagement failure is absent. Final verdict: SENSITIVE. Pathway Agent dissent retained in output. Confidence: 0.87 (decayed from 1.0 by one flip).}
\label{fig:debate-trace}
\end{figure}

% ----------------------------------------------------------
\subsection{(e) Research Design and Methods}
% ----------------------------------------------------------

\subsubsection{Overall Design}

A decentralized multi-agent system where four specialist LLM agents each query a
modality-specific MCP server and produce an EvidencePack, then debate via the
axiom-governed protocol to reach a consensus prediction. The design is motivated by published
tumour board practice \cite{ref12,ref18}, the demonstrated failure of monolithic LLMs on
multi-omics tasks \cite{ref3,ref4}, and the established feasibility of MCP in bioinformatics
\cite{ref15,ref16}. The system is evaluated against two baselines (monolithic LLM and Random
Forest) and through six pre-specified ablations, using both classical predictive metrics and novel
reasoning quality metrics as primary outcomes.

% [IMAGE: Five-tier axiom hierarchy diagram showing T1 Structural Primacy at the apex, descending through T2 Transcriptional Gate, T3 Pathway Bypass, T4 Pharmacological Prior, to T5 Statistical Consensus at the base, with arrows indicating override direction]
\begin{figure}[h]
\centering
\includegraphics[width=0.85\textwidth]{figures/placeholder}
\caption{The five-tier Biological Hierarchy of Truth. T1 (Structural Primacy) governs when DNA-level structural binding failure is detected; T2 (Transcriptional Gate) governs when target gene expression is silenced below functional threshold; T3 (Pathway Bypass) governs when an alternative signalling route is confirmed active; T4 (Pharmacological Prior) governs when historical IC50 data is available for the specific cell line--drug pair; T5 (Statistical Consensus) applies when no higher-tier conflict resolution is determinable. Lower-tier numbers override higher-tier numbers.}
\label{fig:axiom-hierarchy}
\end{figure}

\subsubsection{Datasets}

\textbf{Development set:} 40 (cell line, drug) pairs from GDSC2 \cite{ref20}, stratified across
five axiom tiers (eight cases per tier), with genomic and expression data from CCLE/DepMap
\cite{ref19}. Sensitivity labels derived from z-score normalised IC50: z $< -0.5$ $\to$
SENSITIVE; z $> 0.5$ $\to$ RESISTANT; $|$z$| \leq 0.5$ $\to$ UNCERTAIN. Stratification ensures
that the evaluation exercises every component of the axiom hierarchy; a dataset composed
entirely of T5 cases would test the statistical consensus mechanism but not structural primacy.

\textbf{Validation set:} Held-out CTRPv2 cases, selected to cover the same cancer types and
drug classes present in the development set. CTRPv2 was collected independently of GDSC2,
providing a genuine external validation.

\textbf{Data availability:} All databases are built from publicly available sources. GDSC2 and
CCLE/DepMap data are accessed via the Sanger and Broad Institute public portals. KEGG
pathways are retrieved via the KEGG REST API. ChEMBL drug mechanism data is retrieved via
the ChEMBL public API.

\subsubsection{Aim 1 Methods: Predictive Accuracy}

\textit{Framework run.} Each of the 40 development cases is processed through the full pipeline:
four agents independently query their MCP servers and produce EvidencePacks; ConflictDetector
identifies disagreements; AxiomResolver applies the Hierarchy of Truth; DebateEngine
orchestrates up to three rounds; ConsensusResult is produced with full debate trace. All LLM
calls are cached in a SQLite response cache to enable exact replication of results.

\textit{Multi-model comparison.} The full framework is run independently with five LLM
backends via NVIDIA NIM API. Results are compared across models on all primary metrics to
assess sensitivity of framework performance to underlying LLM quality. The best-performing
model is selected for the ablation study of Aim 2.

\textit{Baselines.}
\begin{itemize}
\item \textit{Monolithic LLM baseline:} Identical context (all four data modalities) provided to
a single LLM in one prompt, without debate or axiom resolution. The monolithic baseline tests
whether the agent architecture adds value beyond simply providing a well-contextualised prompt.
\item \textit{Random Forest baseline:} Trained on GDSC2 features (mutation presence/absence
vectors, expression z-scores, historical IC50 values) using scikit-learn with 5-fold cross-validation.
The Random Forest baseline provides a classical ML anchor point.
\end{itemize}

\textit{Primary metrics:} AUROC and AUPRC (primary); Spearman $\rho$ between predicted
confidence and z-score magnitude (assesses calibration); Cohen's $\kappa$ for label agreement
(assesses beyond-chance performance).

\textit{Potential challenge.} Missing target gene expression data for some cell line/drug
combinations in CCLE. Alternative approach: when direct target expression is absent, broaden
gene lookup to include confirmed pathway members upstream and downstream of the drug target.
This is biologically principled---pathway-level expression is relevant even when the direct target
lacks measurement---and is pre-specified to prevent post-hoc rationalisation.

\subsubsection{Aim 2 Methods: Ablation Study}

Six pre-specified ablations on the same 40 development cases:

\begin{enumerate}
\item \textit{No axiom hierarchy:} Conflict resolved by confidence-weighted majority vote
regardless of evidence tier. Tests whether the axiom ordering adds value over simple confidence
weighting.

\item \textit{Single agent:} All four modalities provided to one agent within the debate wrapper,
eliminating domain isolation. Tests whether specialised agents add value over a single generalist.

\item \textit{Free-text tools:} MCP returns replaced with unstructured natural language
summaries of the same information. Tests whether typed JSON returns add structural anti-hallucination
value.

\item \textit{Randomised tier order:} Axiom hierarchy tier assignments shuffled randomly per
run. Tests whether the specific biological ordering of the hierarchy (T1 $\to$ T5) adds value
over any fixed ordering.

\item \textit{No debate rounds:} Single-round forced resolution after first conflict detection,
eliminating the multi-round DebateEngine. Tests whether iterative argumentation adds value.

\item \textit{Majority vote resolution:} Confidence-weighted vote replaces axiom resolution
throughout. Tests the full system against the standard MAD baseline.
\end{enumerate}

\textit{Reasoning quality metrics.} These metrics are novel contributions to oncology AI evaluation:

\textit{Hallucination rate} is defined as the fraction of biomarkers cited in agent findings that
are absent from the MCP database response for that case. Computed from the structured
EvidencePack: every biomarker field is cross-referenced against the MCP server log for that
agent's queries. A hallucination is a citation of a mutation, expression value, or IC50 that was
not returned by the MCP server.

\textit{Biological faithfulness score} is defined as the fraction of axiom invocations that are
consistent with established molecular causal relationships. Computed against the axiom hierarchy
itself as the reference standard: a T1 override is biologically faithful if and only if the genomics
MCP server returned evidence of a direct structural binding-site alteration; a T3 override is
faithful if and only if the pathways server returned evidence of an active bypass route. The
axiom hierarchy thus serves dual purpose as both the resolution mechanism and the biological
faithfulness reference standard.

\textit{Potential difficulties and mitigations.}

Hallucination rate computation requires automated biomarker extraction from agent findings.
Addressed by requiring structured JSON EvidencePacks with explicit, typed biomarker fields---a
free-text findings field would make automated extraction unreliable. This requirement is enforced
at the MCP protocol layer.

Biological faithfulness scoring requires that axiom invocations be loggable with sufficient
detail to verify against the reference standard. Addressed by the DebateEngine's structured
logging, which records the axiom tier invoked, the evidence cited, and the MCP query ID that
produced that evidence---providing a complete chain of custody for every conflict resolution.

% [IMAGE: Debate protocol state machine diagram showing states: Independent Analysis, EvidencePack Submission, Conflict Detection, Axiom Resolution, Debate Round (up to 3), Consensus or Forced Resolution, Output with Debate Trace]
\begin{figure}[h]
\centering
\includegraphics[width=0.85\textwidth]{figures/placeholder}
\caption{Debate protocol state machine. Agents enter Independent Analysis in parallel, producing EvidencePacks without visibility of other agents' outputs. After all four EvidencePacks are submitted, ConflictDetector transitions to Axiom Resolution if disagreement is detected. DebateEngine orchestrates up to three rounds of structured argumentation. If consensus is reached, the system outputs a ConsensusResult. If three rounds are exhausted without consensus, Forced Resolution applies the highest available axiom tier, and all dissenting positions are retained in the debate trace.}
\label{fig:state-machine}
\end{figure}

\subsubsection{Data Sharing}

All code, curated databases, evaluation cases, and debate traces will be released as open source
on GitHub upon publication. The LLM response cache (SQLite) will be shared to enable exact
replication of results without incurring API costs. All databases are derived from publicly
available sources and will be redistributed under terms consistent with the GDSC2, CCLE, KEGG,
and ChEMBL data use agreements. No patient-identifiable data is involved; all analyses are
performed on cancer cell line data.

% ----------------------------------------------------------
\subsection{(f) Research Timeline}
% ----------------------------------------------------------

% [IMAGE: Gantt chart showing project phases from June 1 through July 31, with milestones for Data Foundation (complete), MCP Servers (complete), Specialized Agents (complete), Consensus Protocol (complete), Evaluation Pipeline (in progress), Live Results and Analysis (in progress), and key dates: July 7 Progress Meet 2, July 16 Midterm Evaluation, July 21 Paper Draft, July 31 Final Presentation]
\begin{figure}[h]
\centering
\includegraphics[width=0.85\textwidth]{figures/placeholder}
\caption{Project Gantt chart covering June 1 through July 31. Completed phases are shown in solid fill; in-progress phases are shown with diagonal hatching. Key milestones: July 7 (Progress Meet 2 --- present preliminary results), July 16 (Midterm Evaluation), July 21 (Paper Draft submission), July 31 (Ready for Final Presentation).}
\label{fig:gantt}
\end{figure}

\begin{center}
\begin{tabular}{lll}
\toprule
\textbf{Phase} & \textbf{Dates} & \textbf{Status} \\
\midrule
Data Foundation & June 1--11 & Complete \\
MCP Servers & June 10--21 & Complete \\
Specialised Agents & June 10--21 & Complete \\
Consensus Protocol & June 15--July 1 & Complete \\
Evaluation Pipeline & June 20--July 31 & In Progress \\
Live Results and Analysis & July 1--31 & In Progress \\
\midrule
\textbf{Milestone} & \textbf{Date} & \textbf{Deliverable} \\
\midrule
Progress Meet 2 & July 7 & Preliminary results presentation \\
Midterm Evaluation & July 16 & Formal evaluation with supervisor \\
Paper Draft & July 21 & Complete manuscript draft \\
Final Presentation & July 31 & Thesis presentation ready \\
\bottomrule
\end{tabular}
\end{center}

\vspace{1em}

All infrastructure layers (Layers 1--4) are complete. The evaluation pipeline (Layer 5) is
implemented and running; the remaining work consists of completing the full Aim 1 multi-model
comparison, executing all six Aim 2 ablations, and analysing and writing up results. The timeline
is feasible given that the primary computational bottleneck---LLM API latency---is addressed
by the SQLite response cache, which enables re-runs of the evaluation pipeline at near-zero
cost after the initial inference pass.

% ----------------------------------------------------------
\subsection{(g) Budget}
\textit{Note: Project budget is funded by mentor, not Johns Hopkins University.}
% ----------------------------------------------------------

\begin{center}
\begin{tabular}{lp{8cm}r}
\toprule
\textbf{Item} & \textbf{Description} & \textbf{Cost} \\
\midrule
NVIDIA NIM API & Free tier access to five hosted LLM backends (llama-3.3-70b, llama-3.2-11b, llama-3.1-8b, nemotron-49b, llama-3.1-70b) & \$0 \\
GDSC2 data access & Publicly available pharmacogenomic dataset (Sanger Institute) & \$0 \\
CCLE/DepMap data access & Publicly available cancer cell line dataset (Broad Institute) & \$0 \\
KEGG API access & Publicly available pathway database & \$0 \\
ChEMBL API access & Publicly available drug mechanism database & \$0 \\
Local compute & MacBook Pro (existing hardware, no rental required) & \$0 \\
Cloud backup & Google Drive (existing subscription) & \$0 \\
Software & Python, FastMCP, SQLite, scikit-learn, matplotlib (all open-source) & \$0 \\
\midrule
\textbf{Total} & & \textbf{\$0} \\
\bottomrule
\end{tabular}
\end{center}

% ----------------------------------------------------------
\subsection{(h) Budget Justification}
% ----------------------------------------------------------

\textbf{Environment and Resources.}
All work is conducted at Johns Hopkins University, which provides institutional support for
biomedical AI research, library access, and high-performance computing resources. The project
requires no wet-lab materials, reagents, or specialised equipment. All computational work runs
on the student's local development environment with cloud-based LLM inference via the NVIDIA
NIM API. The environment is fully reproducible via the project's \texttt{pyproject.toml}
dependency specification, ensuring portability and replicability.

\textbf{Justification of zero-cost budget.}
All computational resources are free or already available. GDSC2 \cite{ref20}, CCLE/DepMap
\cite{ref19}, KEGG \cite{ref16}, and ChEMBL are publicly available research databases provided
at no cost for academic use. NVIDIA NIM provides free API access to large language models
sufficient for the scale of this evaluation (40 development cases, six ablations). All software
(Python, FastMCP, SQLite, scikit-learn, matplotlib) is open-source and freely available.

The student's existing MacBook Pro provides sufficient compute for all data processing and
framework orchestration tasks. The SQLite response cache eliminates redundant API calls,
ensuring that the evaluation pipeline can be run repeatedly---for debugging, analysis, and
paper writing---without incurring additional API cost after the initial inference pass. No wet-lab
materials, reagents, specialized laboratory equipment, or cloud compute rental are required for
any component of this research.

% ============================================================
%  (i) REFERENCES
% ============================================================
\newpage
\subsection*{(i) References}
\begin{thebibliography}{99}

\bibitem{ref1}
Hein ZM, Guruparan D, Okunsai B, Che Mohd Nassir CMN, Ramli MDC, Kumar S.
AI and Machine Learning in Biology: From Genes to Proteins.
\textit{Biology} 2025;14:1453.
\url{https://doi.org/10.3390/biology14101453}

\bibitem{ref2}
El-Sayed AF, Darwish SM, Labib NM, et al.
Generative deep learning in bioinformatics: a microservices-driven paradigm for big data environment.
\textit{J Big Data} 2026;13:7.
\url{https://doi.org/10.1186/s40537-025-01320-5}

\bibitem{ref3}
Zhou S, Xu Z, Zhang M, et al.
Large language models for disease diagnosis: a scoping review.
\textit{npj Artif Intell} 2025;1:9.
\url{https://doi.org/10.1038/s44387-025-00011-z}

\bibitem{ref4}
Meng X, Yan X, Zhang K, et al.
The application of large language models in medicine: a scoping review.
\textit{iScience} 2024;27(5):109713.
doi:10.1016/j.isci.2024.109713.

\bibitem{ref5}
Williams CYK, Miao BY, Kornblith AE, et al.
Evaluating the use of large language models to provide clinical recommendations in the Emergency Department.
\textit{Nat Commun} 2024;15:8236.
\url{https://doi.org/10.1038/s41467-024-52415-1}

\bibitem{ref6}
Yang X, Li T, Su Q, et al.
Application of large language models in disease diagnosis and treatment.
\textit{Chin Med J (Engl)} 2025;138(2):130--142.
doi:10.1097/CM9.0000000000003456.

\bibitem{ref7}
Liu X, Wei K, Lin E, et al.
Multi-omics advances in understanding cancer drug resistance.
\textit{Biomed Technol} 2025;12:100119.
doi:10.1016/j.bmt.2025.100119.

\bibitem{ref8}
Mao Y, Shangguan D, Huang Q, et al.
Emerging artificial intelligence-driven precision therapies in tumor drug resistance: recent advances, opportunities, and challenges.
\textit{Mol Cancer} 2025;24(1):123.
doi:10.1186/s12943-025-02321-x.

\bibitem{ref9}
Saranya KR, Vimina ER.
DRN-CDR: A cancer drug response prediction model using multi-omics and drug features.
\textit{Comput Biol Chem} 2024;112:108175.
doi:10.1016/j.compbiolchem.2024.108175.

\bibitem{ref10}
Zhao H, Zhang X, Zhao Q, Li Y, Wang J.
MSDRP: a deep learning model based on multisource data for predicting drug response.
\textit{Bioinformatics} 2023;39(9):btad514.
\url{https://doi.org/10.1093/bioinformatics/btad514}

\bibitem{ref11}
Liu Q, Hu Z, Huang T, et al.
EvoMDT: a self-evolving multi-agent system for structured clinical decision-making in multi-cancer.
\textit{npj Digit Med} 2026;9:124.
\url{https://doi.org/10.1038/s41746-025-02304-8}

\bibitem{ref12}
Han X, Gao X, Qu X, Yu Z.
Multi-Agent Medical Decision Consensus Matrix System: An Intelligent Collaborative Framework for Oncology MDT Consultations.
arXiv:2512.14321. 2025.

\bibitem{ref13}
Kim Y, Jeong H, Chen S, et al.
Medical Hallucination in Foundation Models and Their Impact on Healthcare.
arXiv:2503.05777. 2025.
doi:10.1101/2025.02.28.25323115.

\bibitem{ref14}
Nishisako S, Higashi T, Wakao F.
Reducing Hallucinations and Trade-Offs in Responses in Generative AI Chatbots for Cancer Information: Development and Evaluation Study.
\textit{JMIR Cancer} 2025;11:e70176.
doi:10.2196/70176.

\bibitem{ref15}
Flotho M, Diks IF, Flotho P, Molano LG, Hirsch P, Keller A.
MCPmed: a call for Model Context Protocol-enabled bioinformatics web services for LLM-driven discovery.
\textit{Brief Bioinform} 2026;27(1):bbag076.
doi:10.1093/bib/bbag076.

\bibitem{ref16}
Widjaja F, Chen Z, Zhou J.
BioinfoMCP: A Unified Platform Enabling MCP Interfaces in Agentic Bioinformatics.
arXiv:2510.02139. 2025.
doi:10.48550/arXiv.2510.02139.

\bibitem{ref17}
Mehandru N, Hall AK, Melnichenko O, et al.
BioAgents: Bridging the gap in bioinformatics analysis with multi-agent systems.
\textit{Sci Rep} 2025;15:39036.
\url{https://doi.org/10.1038/s41598-025-25919-z}

\bibitem{ref18}
Lorenzo L, Montana-Mendez M, Figueiras S, Boubeta M, Bernardo-Casti\~{n}eira C.
Evaluation of Oncotimia: An LLM based system for supporting tumour boards.
arXiv:2601.19899. 2026.
doi:10.48550/arXiv.2601.19899.

\bibitem{ref19}
Barretina J, Caponigro G, Stransky N, et al.
The Cancer Cell Line Encyclopedia enables predictive modelling of anticancer drug sensitivity.
\textit{Nature} 2012;483:603--607.

\bibitem{ref20}
Iorio F, Knijnenburg TA, Vis DJ, et al.
A Landscape of Pharmacogenomic Interactions in Cancer.
\textit{Cell} 2016;166(3):740--754.

\bibitem{ref21}
Turpin M, Michael J, Perez E, Bowman SR.
Language Models Don't Always Say What They Think: Unfaithful Explanations in Chain-of-Thought Prompting.
In: \textit{Advances in Neural Information Processing Systems (NeurIPS)}. 2023.

\bibitem{ref22}
Housman G, Byler S, Heerboth S, et al.
Drug resistance in cancer: an overview.
\textit{Cancers} 2014;6(3):1769--1792.

\bibitem{ref23}
Proietti I, Skroza N, Bernardini N, et al.
Mechanisms of acquired BRAF inhibitor resistance in melanoma: a systematic review.
\textit{Cancers} 2020;12(10):2801.

\bibitem{ref24}
Rotow J, Bivona TG.
Understanding and targeting resistance mechanisms in NSCLC.
\textit{Nat Rev Cancer} 2017;17(11):637--658.

\bibitem{ref25}
Carli F, Di Chiaro P, Morelli M, et al.
Learning and actioning general principles of cancer cell drug sensitivity.
\textit{Nat Commun} 2025;16:1654.
\url{https://doi.org/10.1038/s41467-025-56827-5}

\bibitem{ref26}
Parca L, Petrizzelli M, Conte F, Fiscon G, Paolini A, Paci P.
CellHit: predicting patient drug responses via cancer cell lines.
\textit{Nucleic Acids Res} 2025. (in press)

% -- Additional references from the full PDF extraction --

\bibitem{ref27}
Galitsky BA.
LLM-Based Personalized Recommendations in Health.
Preprints 2024, 2024021709.
\url{https://doi.org/10.20944/preprints202402.1709.v1}

\bibitem{ref28}
Lin C, Kuo CF.
Roles and potential of Large language models in healthcare: A comprehensive review.
\textit{Biomed J} 2025;48.

\bibitem{ref29}
Wu X, Sun M, Li L, Wang J, Wan S.
Leveraging Artificial Intelligence and Large Language Models for Cancer Immunotherapy.
\textit{Advanced Science} 2026:e21936.
\url{https://doi.org/10.1002/advs.202521936}

\bibitem{ref30}
Su Y, Liu J, Li J, et al.
Applications of artificial intelligence and large language models in cancer immunotherapy.
\textit{Current Proteomics} 2025;22(5):100057.
\url{https://doi.org/10.1016/j.curpro.2025.100057}

\bibitem{ref31}
Lai YJ, Wang LJ, Yasaka TM, et al.
Inferring Drug-Gene Relationships in Cancer Using Literature-Augmented Large Language Models.
\textit{Cancer Res Commun} 2025;5(4):706--718.
doi:10.1158/2767-9764.CRC-25-0030.

\bibitem{ref32}
Chen D, Parsa R, Swanson K, et al.
Large language models in oncology: a review.
\textit{BMJ Oncol} 2025;4(1):e000759.
doi:10.1136/bmjonc-2025-000759.

\bibitem{ref33}
Li A, Zhang J, Li L, et al.
M3MAD-Bench: Are Multi-Agent Debates Really Effective Across Domains and Modalities?
arXiv:2601.02854. 2026.
doi:10.48550/arXiv.2601.02854.

\bibitem{ref34}
Yang T, Xiao Y, Bao Z, Hao J, Peng J.
The rise and potential opportunities of large language model agents in bioinformatics and biomedicine.
\textit{Brief Bioinform} 2025;26(6):bbaf601.
doi:10.1093/bib/bbaf601.

\bibitem{ref35}
Zhou J, Jiang J, Han Z, Wang Z, Gao X.
Streamline automated biomedical discoveries with agentic bioinformatics.
\textit{Brief Bioinform} 2025;26(5):bbaf505.
doi:10.1093/bib/bbaf505.

\bibitem{ref36}
Martinek V, Gariboldi A, Tzimotoudis D, et al.
Agentomics-ML: Autonomous Machine Learning Experimentation Agent for Genomic and Transcriptomic Data.
arXiv:2506.05542. 2025.
doi:10.48550/arXiv.2506.05542.

\bibitem{ref37}
Zack M, Stupichev DN, Moore AJ, et al.
Artificial Intelligence and Multi-Omics in Pharmacogenomics: A New Era of Precision Medicine.
\textit{Mayo Clin Proc Digit Health} 2025;3(3):100246.
doi:10.1016/j.mcpdig.2025.100246.

\bibitem{ref38}
He H, et al.
Biology-Instructions: A Dataset and Benchmark for Multi-Omics Sequence Understanding Capability of Large Language Models.
In: \textit{Findings of the Association for Computational Linguistics: EMNLP 2025}, pages 17984--18016.

\bibitem{ref39}
Qu S, Ding N, Xie L, et al.
Automating Exploratory Multiomics Research via Language Models.
arXiv:2506.07591. 2025.
doi:10.48550/arXiv.2506.07591.

\bibitem{ref40}
Ali S, Qadri YA, Ahmad K, et al.
Large Language Models in Genomics---A Perspective on Personalized Medicine.
\textit{Bioengineering (Basel)} 2025;12(5):440.
doi:10.3390/bioengineering12050440.

\bibitem{ref41}
Dip SA, Zafor A, Paul B, et al.
LLM4Cell: A Survey of Large Language and Agentic Models for Single-Cell Biology.
arXiv:2510.07793. 2025.
doi:10.48550/arXiv.2510.07793.

\bibitem{ref42}
Vavekanand R.
A Comprehensive Review of Multimodal Large Language Models for Medical Imaging and Omics Data.
\textit{Arch Comput Methods Eng} 2026.
doi:10.1007/s11831-026-10504-y.

\bibitem{ref43}
Yang X, Shashidhar KC, Zhang M, et al.
PromptBio: A Multi-Agent AI Platform for Bioinformatics Data Analysis.
bioRxiv:2025.07.05.663295. 2025.
doi:10.1101/2025.07.05.663295.

\bibitem{ref44}
Du P, Fan R, Zhang N, Wu C, Zhang Y.
Advances in Integrated Multi-omics Analysis for Drug-Target Identification.
\textit{Biomolecules} 2024;14(6):692.
doi:10.3390/biom14060692.

\bibitem{ref45}
Chen HO, Cui YC, Lin PC, Chiang JH.
An Innovative Multi-Omics Model Integrating Latent Alignment and Attention Mechanism for Drug Response Prediction.
\textit{J Pers Med} 2024;14:694.
\url{https://doi.org/10.3390/jpm14070694}

\bibitem{ref46}
Jiang W, Ye W, Tan X, et al.
Network-based multi-omics integrative analysis methods in drug discovery: a systematic review.
\textit{BioData Mining} 2025;18:27.
\url{https://doi.org/10.1186/s13040-025-00442-z}

\bibitem{ref47}
Almutiri T, Alomar K, Alganmi N.
Predicting Drug Response on Multi-Omics Data Using a Hybrid of Bayesian Ridge Regression with Deep Forest.
\textit{Int J Adv Comput Sci Appl} 2023;14.
doi:10.14569/IJACSA.2023.0140550.

\bibitem{ref48}
Liu XY, Mei XY.
Prediction of drug sensitivity based on multi-omics data using deep learning and similarity network fusion approaches.
\textit{Front Bioeng Biotechnol} 2023;11:1156372.
doi:10.3389/fbioe.2023.1156372.

\bibitem{ref49}
Vasilev K, Misrahi A, Jain E, et al.
MTBBench: A Multimodal Sequential Clinical Decision-Making Benchmark in Oncology.
arXiv:2511.20490. 2025.
doi:10.48550/arXiv.2511.20490.

\bibitem{ref50}
Zhang R, Wang J, Yang S, et al.
Multi-Agent Intelligence for Multidisciplinary Decision-Making in Gastrointestinal Oncology.
arXiv:2512.08674. 2025.
doi:10.48550/arXiv.2512.08674.

\bibitem{ref51}
Codella N, Preston S, Qiu H, et al.
Demo: Healthcare Agent Orchestrator (HAO) for Patient Summarization in Molecular Tumor Boards.
arXiv:2509.06602. 2025.
doi:10.48550/arXiv.2509.06602.

\bibitem{ref52}
Kaesberg LB, Becker J, Wahle JP, Ruas T, Gipp B.
Voting or Consensus? Decision-Making in Multi-Agent Debate.
In: \textit{Findings of the Association for Computational Linguistics: ACL 2025}, pages 11640--11671, Vienna, Austria.

\bibitem{ref53}
Spieser J, Balapour A, Meller J, Patra KC, Shamsaei B.
A Review of Multi-Agent AI Systems for Biological and Clinical Data Analysis.
\textit{Methods Protoc} 2026;9(2):33.
\url{https://doi.org/10.3390/mps9020033}

\bibitem{ref54}
Liu Y, Yao H, Liu W, He A, Zhang Y.
Heterogeneous Consensus-Progressive Reasoning for Efficient Multi-Agent Debate.
arXiv:2604.09679. 2026.

\bibitem{ref55}
Liu T, Wang X, Huang W, et al.
GroupDebate: Enhancing the Efficiency of Multi-Agent Debate Using Group Discussion.
arXiv:2409.14051. 2024.

\bibitem{ref56}
Cui Y, Fu H, Zhang H, Wang L, Zuo C.
Free-MAD: Consensus-Free Multi-Agent Debate.
arXiv:2509.11035. 2025.

\bibitem{ref57}
Ngoc Nguyen O, Amin D, Bennett J, et al.
GP or ChatGPT? Ability of large language models to support general practitioners when prescribing antibiotics.
\textit{J Antimicrob Chemother} 2025;80(5):1324--1330.
doi:10.1093/jac/dkaf077.

\bibitem{ref58}
Testagrose C, Pandey S, Serajian M, et al.
Leveraging large language models to predict antibiotic resistance in \textit{Mycobacterium tuberculosis}.
\textit{Bioinformatics} 2025;41(Supplement\_1):i40--i48.
\url{https://doi.org/10.1093/bioinformatics/btaf232}

\bibitem{ref59}
Yoo H, Sokhansanj B, Brown JR, Rosen G.
Predicting Anti-microbial Resistance using Large Language Models.
arXiv:2401.00642. 2024.

\bibitem{ref60}
De Vito A, Geremia N, Bavaro DF, et al.
Comparing large language models for antibiotic prescribing in different clinical scenarios: which performs better?
\textit{Clin Microbiol Infect} 2025;31(8):1336--1342.
\url{https://doi.org/10.1016/j.cmi.2025.03.002}

\end{thebibliography}

\end{document}
